\chapter{Úvod}
\label{sec:Introduction}
S neustálým rozvojem informačních technologií pronikají počítačové systémy do všech aspektů lidského života, kde se stále častěji setkáváme s reaktivními systémy.
Příkladem těchto systémů mohou být běžné webové služby, nápojové automaty, ale i kritické řídicí systémy v dopravě, zdravotnictví a průmyslu. 
Tyto systémy nejsou navrženy k jednorázovému výpočtu a následnému ukončení, ale k~neustálé interakci se svým okolím. 
Vzhledem k jejich paralelní a nedeterministické povaze je~zaručení jejich bezchybného chování pomocí tradičních způsobů testování prakticky nemožné. 

Pro zaručení korektnosti systémů se proto využívají matematicky podložené metody formální verifikace. 
Jednou ze základních metod v této oblasti je kalkulus komunikujících systémů, který představil britský informatik Robin Milner. 
Tento formální jazyk umožňuje modelovat složité konkurentní systémy pomocí propojení jednoduchých nezávislých procesů, zkoumat jejich chování přes strukturální operační sémantiku a vizualizovat je pomocí ohodnocených přechodových systémů.

Výuka těchto systémů ale často naráží na překážky, a to zejména v prvotních hodinách, kdy se~studenti setkají s formálními matematickými důkazy. 
Existující nástroje jsou totiž často navrženy pro potřeby pokročilé analýzy, nejsou tedy přizpůsobeny pro interaktivitu a výuku úplných základů.
Cílem této diplomové práce je proto navrhnout a implementovat výukovou webovou aplikaci, která studentům usnadní pochopení základních problémů z oblasti formální verifikace systémů. 
Práce si klade za cíl vytvořit dynamickou a uživatelsky přívětivou aplikaci, ve které mohou uživatelé experimentovat s vlastními modely a vizualizovat jejich chování. 
Aplikace zajišťuje kontrolu syntaxe zadaných CCS výrazů, zobrazení jejich syntaktických stromů a poskytuje uživateli prostředí pro~interaktivní a kontrolovanou tvorbu formálních důkazů podle pravidel SOS. 
Vedle toho systém umožňuje automatické generování odpovídajících LTS grafů a simulaci běhu jednotlivých procesů. 
Aby byla aplikace co nejpřístupnější, obsahuje také sadu připravených ukázkových vstupů.

Text práce je rozdělen do pěti na sebe navazujících kapitol. 
V kapitole \ref{sec:reaktivni_systemy} se probírá teoretický základ reaktivních systémů a detailně popisuje formalismus kalkulu komunikujících systémů, ohodnocených přechodových systémů a strukturální operační sémantiky. 
Kapitola \ref{sec:analyza_pozadavku} analyzuje existující řešení a definuje funkční i nefunkční požadavky na vyvíjenou aplikaci. 
V kapitole \ref{sec:technologie} jsou představeny použité moderní webové technologie a knihovny, jako jsou React.js, Peggy.js a Cytoscape.js. 
Kapitola \ref{sec:navrh} se věnuje celkovému architektonickému návrhu řešení a uživatelského rozhraní.
V neposlední řadě popisuje kapitola \ref{sec:implementace} samotnou implementaci výpočetní logiky, datových struktur, algoritmů a~způsob nasazení aplikace.

Výsledná verze aplikace je za pomoci nástroje Github Pages \cite{github_pages_2026_github} nazasena na veřejné adrese \\ \texttt{https://ccs.koberskyj.cz/}.

Během vývoje této práce byla využívána umělá inteligence, konkrétně model Gemini pro jazykovou korekturu textu a pro asistenci při programování, zejména tedy při ladění aplikace v prostředí React.js.

\endinput